% Figure: two nonlinear input-output characteristics that a single-valued
% function cannot capture. Left: a dead-band-plus-saturation static
% characteristic (still a function). Right: a hysteresis loop, where the
% output depends on the direction of travel and so is NOT a function of
% the instantaneous input.
\begin{figure}[htbp]
\centering
\resizebox{\linewidth}{!}{%
\begin{tikzpicture}[scale=1.0,
  ax/.style={->, draw=engblack, thick},
  tick/.style={draw=engblack!60, thin},
  curve/.style={draw=engblue, very thick},
  up/.style={draw=engblue, very thick},
  down/.style={draw=engred, very thick},
  ann/.style={font=\footnotesize, align=center},
]

% --- Left: dead-band plus saturation (single-valued) ---
\begin{scope}[xshift=0cm]
  \draw[ax] (-3.2, 0) -- (3.2, 0) node[below right, ann] {input $u$};
  \draw[ax] (0, -2.2) -- (0, 2.2) node[left, ann] {output $y$};
  % dead-band [-0.8, 0.8]; linear region to +/-2.0; saturate at +/-1.6
  \draw[curve] (-3.0, -1.6) -- (-2.0, -1.6); % low saturation
  \draw[curve] (-2.0, -1.6) -- (-0.8, 0);    % rising through dead-band edge
  \draw[curve] (-0.8, 0) -- (0.8, 0);        % dead band (zero output)
  \draw[curve] (0.8, 0) -- (2.0, 1.6);       % rising
  \draw[curve] (2.0, 1.6) -- (3.0, 1.6);     % high saturation
  \draw[tick] (-0.8, -0.05) -- (-0.8, 0.05);
  \draw[tick] (0.8, -0.05) -- (0.8, 0.05);
  \node[ann, below, yshift=-2pt] at (0, -0.15) {dead band};
  \node[ann, anchor=south west, fill=white, inner sep=1pt] at (2.0, 1.6)
    {saturation};
  \node[ann, anchor=north] at (0, -2.5)
    {(a) dead band $+$ saturation\\(a function: one $y$ per $u$)};
\end{scope}

% --- Right: hysteresis loop (not single-valued) ---
\begin{scope}[xshift=8.5cm]
  \draw[ax] (-3.2, 0) -- (3.2, 0) node[below right, ann] {input $u$};
  \draw[ax] (0, -2.2) -- (0, 2.2) node[left, ann] {output $y$};
  % rising branch shifted right; falling branch shifted left
  \draw[up] (-3.0, -1.6) -- (-2.0, -1.6) -- (1.2, 1.6) -- (3.0, 1.6);
  \draw[down] (3.0, 1.6) -- (2.0, 1.6) -- (-1.2, -1.6) -- (-3.0, -1.6);
  % arrows on branches
  \draw[up, -{Stealth[length=4pt]}] (-0.6, 0.0) -- (-0.4, 0.1);
  \draw[down, -{Stealth[length=4pt]}] (0.6, 0.0) -- (0.4, -0.1);
  \node[ann, text=engblue, anchor=south east] at (1.0, 1.2)
    {increasing $u$};
  \node[ann, text=engred, anchor=north west] at (-1.0, -1.2)
    {decreasing $u$};
  \node[ann, anchor=north] at (0, -2.5)
    {(b) hysteresis loop\\(not a function: $y$ depends on history)};
\end{scope}

\end{tikzpicture}%
}%
\caption[Dead band, saturation, and hysteresis]{Two nonlinear
input-output characteristics. Panel (a): a dead band (a window of
input near zero that produces no output, common in valves, gear
trains with backlash, and threshold comparators) followed by a linear
region and saturation at the rails. The characteristic is still a
single-valued function, one output for each input, so the chapter's
machinery applies directly. Panel (b): a hysteresis loop, where the
output on the rising branch differs from the output on the falling
branch at the same input value. A magnetic material's $B$-$H$ curve,
a thermostat with a switching differential, and a mechanism with
static friction all show this shape. Hysteresis is not a function of
the instantaneous input: the output depends on the direction and
history of travel, so it requires a state variable beyond the input
to predict. Recognising that a characteristic has crossed from
function to history-dependent map is the first step in modelling it
correctly.}
\label{fig:vol02:ch02:hysteresis-deadband}
\end{figure}
